% Part: normal-modal-logic
% Chapter: syntax-and-semantics
% Section: relational-models

\documentclass[../../../include/open-logic-section]{subfiles}

\begin{document}

\olfileid{nml}{syn}{rel}

\olsection{مدل‌های رابطه‌ای}

مفهوم پایهٔ معناشناسی برای منطق‌های موجهات عادی، مفهوم
\emph{مدل رابطه‌ای} است. چنین مدلی از مجموعه‌ای از جهان‌ها تشکیل می‌شود
که با یک «رابطهٔ دسترسی» دوتایی به یکدیگر مرتبط‌اند، به‌همراه یک
تخصیص که تعیین می‌کند کدام !!{propositional variable} در کدام جهان
«صادق» به شمار می‌آید.

\begin{defn}
  یک \emph{مدل} برای زبان پایهٔ موجهات سه‌تاییِ~$\mModel{M}
  = \tuple{W, R, V}$ است، که در آن
  \begin{enumerate}
  \item $W$ مجموعه‌ای ناتهی از «جهان‌ها» است؛
  \item $R$ یک رابطهٔ دسترسی دوتایی روی~$W$ است؛ و
  \item $V$ تابعی است که به هر !!{propositional
    variable}~$p$ مجموعه‌ای از جهان‌های ممکن، یعنی~$V(p)$، نسبت می‌دهد.
  \end{enumerate}
  هرگاه~$Rww'$ برقرار باشد، می‌گوییم~$w'$ \emph{از~$w$ دسترس‌پذیر}
  است. هرگاه~$w \in V(p)$، می‌گوییم~$p$ \emph{در~$w$ صادق} است.
\end{defn}

مزیت بزرگ معناشناسی رابطه‌ای این است که مدل‌ها را می‌توان با نمودارهای
ساده‌ای مانند نمودار~\olref{fig:simple} نمایش داد. جهان‌ها با گره‌ها
نمایش داده می‌شوند، و جهان~$w'$ دقیقاً هنگامی از~$w$ دسترس‌پذیر است که
پیکانی از~$w$ به~$w'$ وجود داشته باشد. افزون‌براین، هرگاه
$w \in V(p)$ باشد، گره (جهان) را با~$\mTrue{p}$، و در غیر این صورت
با~\mFalse{p} برچسب می‌زنیم. \olref{fig:simple} مدلی را نمایش می‌دهد
که در آن~$W = \{w_1, w_2, w_3\}$،
$R = \{\tuple{w_1, w_2},\tuple{w_1, w_3}\}$،
$V(p) = \{w_1, w_2\}$ و~$V(q) = \{w_2\}$ است.

\begin{figure}
  \begin{center}
    \begin{tikzpicture}[modal]
      \node[world] (w1) [label={[align=right]right:\mTrue{p}\\ \mFalse{q}}]{$w_1$}; 
      \node[world] (w2) [label={[align=right]right:\mTrue{p}\\ \mTrue{q}},
        above right=of w1]{$w_2$}; 
      \node[world] (w3) [label={[align=right]right:\mFalse{p}\\ \mFalse{q}},
        below right=of w1] {$w_3$};
      \draw[->] (w1) to (w2);
      \draw[->] (w1) to (w3);
    \end{tikzpicture}
  \end{center}
  \caption{یک مدل ساده.}
  \ollabel{fig:simple}
\end{figure}

\end{document}
